\chapter[Introducción]{Introducción}

El presente documento describe el desarrollo de \emph{OnCuidar}, una aplicación
móvil que apoya a los cuidadores de niños, niñas y adolescentes con cáncer
durante el seguimiento domiciliario que sigue al alta hospitalaria. El trabajo
recorre el proceso completo, desde el estudio de la situación actual hasta los
planes de capacitación e implantación del prototipo.

El documento se organiza de la siguiente manera. El capítulo~2 presenta la
institución, el área de estudio y la problemática que motiva el proyecto. El
capítulo~3 define los objetivos y el ambiente de ingeniería de software. El
capítulo~4 especifica los requerimientos del sistema según el estándar
IEEE~830 y el capítulo~5 evalúa su factibilidad técnica, operativa y económica.

Los capítulos~6 y~7 contienen el análisis y el diseño de la solución, con el
proceso de negocio futuro, los casos de uso, el modelo de datos, la arquitectura
y la interfaz. El capítulo~8 documenta las pruebas según el estándar IEEE~829.
Los capítulos~9 y~10 plantean la capacitación y la puesta en marcha, el
capítulo~11 resume el esfuerzo invertido y el capítulo~12 reúne las conclusiones.

Al final se incluyen la bibliografía y los anexos de planificación, resultados
por sprint, manual de usuario, especificación de pruebas, diccionario de datos,
trazabilidad, consideraciones éticas y especificación de los casos de uso
restantes.
